\documentclass{my_eth_beamer}

\usepackage{my_math_eth}

\numberwithin{equation}{section}

% -----------------------------------------------------------------------------
% this document
\title{Introduction to Survey Sampling}
\newcommand{\TITLE}{Variance estimation}
\newcommand{\NR}{7}
\subtitle{Part \NR: \TITLE}
\author[Tobias Schoch]{Tobias Schoch}
\date{2024}

% -----------------------------------------------------------------------------
\begin{document}

\section{\TITLE}\label{sec:variance}
\setcounter{section}{\NR}
\setcounter{Bsp}{19}
\setcounter{Def}{13}
\setcounter{EinSatz}{1}

\setstretch{1.1}

% -----------------------------------------------------------------------------
\begin{frame}
    \frametitle{Overview}
    \noindent \hspace{0.2em}\textcolor{fhnwblue}{\bfseries \NR~\TITLE}
    \begin{enumerate}[\NR.1]
    \setlength{\itemsep}{4pt}
        \item Taylor series expansion (linearization)
        \item Resampling methods
            \begin{enumerate}[a)]
            \setlength{\itemsep}{4pt}
                \item Bootstrap
                \item Jackknife
            \end{enumerate}
    \end{enumerate}
\end{frame}

% -----------------------------------------------------------------------------
\begin{frame}
    \frametitle{\insertsectionnumber~\insertsection}
    \begin{itemize}
    \setlength{\itemsep}{10pt}
        \item The variance estimators introduced in Chapters 2--6 were derived
            for estimators of population totals and means, but not directly
            apply for estimators of \alert{nonlinear} population
            characteristics.
        \item An important \alert{class} of nonlinear characteristics can be
            written as a \alert{function} of a vector of population totals,
            \begin{equation*}
                \bold t = (t_1, \ldots, t_q)^T, \qquad q \geq 1.
            \end{equation*}
        \item The best-known member of this class is the \alert{population
            ratio} of two totals, $\rat = g(\bold t)$, with $\bold
            t=(\ty,\tx)^T$, given by (provided that $\tx \neq 0$)
            \begin{equation*}
                \rat(\bold t) = \frac{\ty}{\tx} = \frac{\myu}{\mx}.
            \end{equation*}
    \end{itemize}
\end{frame}

% -----------------------------------------------------------------------------
\begin{frame}
    \begin{itemize}
    \setlength{\itemsep}{10pt}
        \item For the remainder of this section, assume an \SRS of size $n$.
        \item The estimator
            \begin{equation*}
            \widehat R = \frac{\widehat t_y}{\widehat t_x}
            \end{equation*}
            is commonly referred to as the \alert{ratio estimator}. It is
            generally \alert{biased} for $\rat$.
        \item The \alert{absolute bias relative to the standard error} is
            bounded above by \citep{hartley_ross_1954}
            \begin{equation*}
                \frac{\big\vert \bias\big(\rats\,\big) \big\vert}{%
                    \se\big(\rats\,\big)} \leq
                    \sqrt{\frac{1}{n} - \frac{1}{N}} \mathsf{CV}_x,
            \end{equation*}
            where $\mathsf{CV}_x = \sigma_x / \mx$ is the coefficient of
            variation.
            \vspace{0.25em}
            \begin{itemize}
            \setlength{\itemsep}{10pt}
                \item[\goto] Thus, the bias is small relative to the
                    estimator's standard error when the sample is sufficiently
                    large and $\mathsf{CV}_x$ is not too large.
            \end{itemize}
    \end{itemize}
\end{frame}

\subsection{Taylor series expansion (linearization)}

\begin{frame}
    \frametitle{\insertsectionnumber.\insertsubsectionnumber~\insertsubsection}
    \begin{itemize}
    \setlength{\itemsep}{12pt}
        \item We restrict attention to the ratio estimator $\rats$ under
            \SRS. The same idea applies to other nonlinear characteristics
            that are differentiable functions of population totals.
        \item Write
            \begin{equation*}
                \rats = g(\mys,\mxs), \qquad g(a,b)=\frac{a}{b},
            \end{equation*}
            where $\mys$ and $\mxs$ are the sample means.
        \item A first-order \alert{Taylor expansion} around
            $(\myu, \mx)$ gives
            \begin{equation}\label{eq:variance_taylor}
                \rats = \rat + \frac{1}{\mx} \big(\mys - \rat\, \mxs\big)
                + \op\big(n^{-1/2}\big).
            \end{equation}
    \end{itemize}
\end{frame}

% -----------------------------------------------------------------------------
\begin{frame}
    \begin{itemize}
    \setlength{\itemsep}{9pt}
        \item Let $r_n$ denote the remainder in the Taylor expansion. Under the
            asymptotic framework introduced in Chapter 2,
            \begin{equation*}
                r_n = \op\big(n^{-1/2}\big)
                \qquad \text{means that} \qquad
                \frac{r_n}{n^{-1/2}} \; \xrightarrow{p} \; 0.
            \end{equation*}
        \item Ignoring the remainder in (\ref{eq:variance_taylor}) gives the
            \alert{linearization}
            \begin{equation*}
                \rats_0 = \rat + \frac{1}{\mx} \big(\mys-\rat\,\mxs\big).
            \end{equation*}
        \item Define the \alert{linearized variable}
            \begin{equation*}
                z_i = y_i - \rat x_i.
            \end{equation*}
            Then
            \begin{equation*}
                \rats_0 - \rat = \frac{\bar z}{\mx},
                \qquad
                \var(\rats) \approx \var(\rats_0) =
                \frac{1}{\mx^2}\var(\bar z).
            \end{equation*}
    \end{itemize}
\end{frame}

% -----------------------------------------------------------------------------
\begin{frame}
    \begin{itemize}
    \setlength{\itemsep}{9pt}
        \item Under \SRS, the linearization yields
            \begin{equation*}
                \var(\rats) \approx \var(\rats_0) =
                \left(1-\frac{n}{N}\right)
                \frac{\sigma_z^2}{n\,\mx^2},
            \end{equation*}
            where $\sigma_z^2$ is the population variance of the linearized
            variable.
        \item Since $\rat$ and $\mx$ are unknown, define the
            \alert{estimated linearized variable}
            \begin{equation*}
                \widehat z_i = y_i - \rats x_i.
            \end{equation*}
        \item A corresponding \alert{variance estimator} is
            \begin{equation*}
                \varest(\rats) =
                \left(1-\frac{n}{N}\right)
                \frac{\widehat{\sigma}_{\,\widehat{z}}^{\,2}}{n\,\mxs^2},
                \qquad \text{where} \qquad
                \widehat{\sigma}_{\,\widehat z}^{\,2} =
                \frac{1}{n-1}
                \sum_{i\in s}\widehat z_i^{\,2}.
            \end{equation*}
    \end{itemize}
\end{frame}

% -----------------------------------------------------------------------------
\begin{frame}
    \begin{itemize}
    \setlength{\itemsep}{10pt}
        \item Under suitable regularity conditions, Taylor linearization can
            also be applied to \alert{complex sampling designs}; see, e.g.,
            \citet[][Chapter 2.1]{fuller_2009}.

        \item A main \alert{disadvantage} is that the linearized variable must
            be derived for each nonlinear estimator separately. This can
            become cumbersome for complex estimands and sampling designs.

        \item \alert{Resampling methods} provide an alternative approach to
            variance estimation that often requires less estimator-specific
            derivation.
    \end{itemize}
\end{frame}

\subsection{Resampling methods}

% -----------------------------------------------------------------------------
\begin{frame}
    \frametitle{\insertsectionnumber.\insertsubsectionnumber~\insertsubsection}
    \begin{itemize}
    \setlength{\itemsep}{9pt}
        \item \alert{Resampling methods} estimate the variance of an estimator
            by repeatedly constructing \alert{replicate samples} from the
            observed sample and recomputing the estimator.
        \item Their main advantage is that they often avoid
            \alert{estimator-specific analytical derivations}, such as those
            required for Taylor linearization.
        \item The price is additional \alert{computational effort}, especially
            for complex sampling designs.
        \item We consider two resampling methods:
            \vspace{0.25em}
            \begin{enumerate}[a)]
            \setlength{\itemsep}{4pt}
                \item Bootstrap
                \item Jackknife
            \end{enumerate}
    \end{itemize}
\end{frame}

% -----------------------------------------------------------------------------
\begin{frame}
    \noindent\textcolor{fhnwblue}{\bfseries \large Bootstrap}

    \vspace{0.5em}

    \begin{itemize}
    \setlength{\itemsep}{12pt}
        \item Let $s$ be an \SRS of size $n$ and let $\widehat{\theta}$ be an
            estimator of the population characteristic $\theta_U$.
        \item Choose the number of \alert{bootstrap replicates} $B$.
        \item For each $b = 1, \ldots, B$:
            \vspace{0.25em}
            \begin{enumerate}
            \setlength{\itemsep}{7pt}
                \item Draw a \alert{bootstrap sample} $s_b^*$ of size $n$
                    from $s$, sampling \alert{with replacement}.
                \item Compute the estimator from $s_b^*$ and denote the result
                    by $\widehat{\theta}_b^*$.
            \end{enumerate}
        \item The resulting values
            $\widehat{\theta}_1^*,\ldots,\widehat{\theta}_B^*$
            approximate the sampling variability of $\widehat{\theta}$.
    \end{itemize}
\end{frame}

% -----------------------------------------------------------------------------
\begin{frame}
    \begin{itemize}
    \setlength{\itemsep}{9pt}
        \item The \alert{bootstrap variance estimator} is the sample variance
            of the replicate estimates:
            \begin{equation*}
                \varest^{\mathrm{boot}\;}(\widehat{\theta}) = \frac{1}{B-1}
                \sum_{b=1}^B
                \left(\widehat{\theta}_b^*
                -\widebar{\theta}^{\,*}\right)^2, \qquad
                \widebar{\theta}^{\,*} =
                \frac{1}{B}\sum_{b=1}^B\widehat{\theta}_b^*.
            \end{equation*}
        \item Under \SRS, the ordinary bootstrap is appropriate when the
            \alert{sampling fraction}
            \begin{equation*}
            f=\frac{n}{N}
            \end{equation*}

            \vspace{-1.5em}

            is small.
        \item If $f$ is not small, sampling without replacement produces less
            variability than the ordinary bootstrap. For smooth estimators,
            this can be accounted for approximately by the \alert{finite
            population correction} \citep[][Chapter 5]{wolter_2007}
            \vspace{-0.5em}
            \begin{equation*}
                (1-f)\, \varest^{\mathrm{boot}}(\widehat{\theta}).
            \end{equation*}
    \end{itemize}
\end{frame}

% -----------------------------------------------------------------------------
\begin{frame}
    \frametitle{\normalsize{Bootstrap for Stratified Sampling}}
    \begin{itemize}
    \setlength{\itemsep}{8pt}
        \item For a stratified design, the bootstrap should
            \alert{preserve the stratification}.
        \item For each bootstrap replicate $b=1,\ldots,B$:
            \vspace{0.25em}
            \begin{enumerate}
            \setlength{\itemsep}{5pt}
                \item Within each stratum $h = 1, \ldots, H$, draw a bootstrap
                    sample $s_{hb}^*$ of size $n_h$ from $s_h$, sampling
                    \alert{with replacement}.
                \item Combine the stratum-specific bootstrap samples and
                    compute the replicate estimate $\widehat{\theta}_b^*$.
            \end{enumerate}
        \item If the stratum-specific sampling fractions $f_h = n_h / N_h$
            are small, the ordinary within-stratum bootstrap provides a
            useful approximation.
        \item If the $f_h$ are not small, a single finite population
            correction is generally \alert{not sufficient}.
    \end{itemize}
\end{frame}

% -----------------------------------------------------------------------------
\begin{frame}
    \frametitle{\normalsize{Rao--Wu Rescaling Bootstrap}}
    \begin{itemize}
    \setlength{\itemsep}{8pt}
        \item If the stratum-specific sampling fractions $f_h$ are not small,
            the ordinary within-stratum bootstrap does not adequately reflect
            sampling without replacement.
        \item The \alert{Rao--Wu rescaling bootstrap} modifies the bootstrap
            procedure to account for this; see \citet{rao_wu_1988}.
        \item Each bootstrap replicate is associated with a set of adjusted
            \alert{replicate weights}. The estimator is recomputed using these
            weights, yielding replicate estimates
            \begin{equation*}
                \widehat{\theta}_1^*,\ldots,\widehat{\theta}_B^*.
            \end{equation*}
        \item The approach can be applied to many nonlinear estimators without
            deriving a separate linearized variable, including estimators
            based on \alert{quantiles}.
    \end{itemize}
\end{frame}

% -----------------------------------------------------------------------------
\begin{frame}
    \frametitle{\large Jackknife}
    \begin{itemize}
    \setlength{\itemsep}{12pt}
        \item Let $s$ be an \SRS of size $n$ from $U$.
        \item We are interested in the variance of an estimator
            $\widehat{\theta}$ of the characteristic $\theta_U$, where
            $\theta_U$ is defined as a \alert{differentiable} function
            $g(\bold t)$ of population totals $\bold t$.
        \item \alert{Differentiability} of $g$ is a \emph{necessary}
            condition \citep{krewski_rao_1981} for consistency of the
            (delete-1) jackknife variance estimator.
        \item In the \alert{delete-1 jackknife}, one observation at a time is
            omitted from $s$, and the estimator is recomputed without that
            observation.
        \item We write $\widehat{\theta}_{(-j)}$ for the estimator
            $\widehat{\theta}$ computed \alert{without} the $j$th observation
            ($j\in s$).
    \end{itemize}
\end{frame}

% -----------------------------------------------------------------------------
\begin{frame}
    \begin{itemize}
    \setlength{\itemsep}{10pt}
        \item For example, let $\widehat{\theta}:=\mys$, then
            \begin{equation*}
                \widehat{\theta}_{(-j)} = \frac{1}{n - 1}
                    \sum_{i\in s\setminus j} y_i.
            \end{equation*}
        \item The delete-1 jackknife \alert{variance estimator} for
            $\widehat{\theta}$ is
            \begin{equation*}
                \varest^{\mathrm{JK}}(\widehat{\theta}\,) = \frac{n}{(n
                - 1)} \sum_{j = 1}^n \left(\widehat{\theta}_{(-j)} -
                \widehat{\theta}\,\right)^2.
            \end{equation*}
        \item The variance estimator is reliable when the sampling fraction $f$
            is \alert{small}.
            \vspace{0.5em}
            \begin{itemize}
            \setlength{\itemsep}{2pt}
                \item[\goto] If $f$ is \alert{not small}, it can be
                    substantially biased. In this case, it is recommended to
                    multiply it by $(1-f)$; see e.g.
                    \citet[][p.~167]{wolter_2007}.
            \end{itemize}
    \end{itemize}
\end{frame}

% -----------------------------------------------------------------------------
\begin{frame}
    \begin{itemize}
    \setlength{\itemsep}{10pt}
        \item The jackknife variance estimator can also be used with
            \alert{complex designs}; see, e.g., \citet[][Chapter 4]{wolter_2007}.
        \item The jackknife can also be generalized by deleting $d>1$
            observations at a time, i.e., the \alert{delete-$d$ jackknife}.
            \vspace{0.5em}
            \begin{itemize}
            \setlength{\itemsep}{2pt}
                \item[\goto] \citet{shao_wu_1989} further showed that, with an
                    appropriate choice of $d$, the delete-$d$ jackknife
                    variance estimator can be consistent for statistics such as
                    the median or a linear combination of quantiles, unlike the
                    delete-1 jackknife.
            \end{itemize}
    \end{itemize}
\end{frame}

% -----------------------------------------------------------------------------
\setcounter{section}{\the\numexpr\value{section}-1\relax}
\begin{frame}
    \frametitle{References}
    {\footnotesize
        \setstretch{1.1}
        \bibliographystyle{fhnw_en}
        \bibliography{master}
    }
\end{frame}

% -----------------------------------------------------------------------------
\lastpage

\end{document}
